% Viva preparation: the questions an examiner is likely to ask, with short answers built on the project's own numbers.
% Numbers are \R{...} keys from ../report/generated/numbers.tex, so they always match the report.
\documentclass[11pt,a4paper]{article}
\usepackage[a4paper,margin=22mm]{geometry}
\usepackage[T1]{fontenc}
\usepackage{lmodern}
\usepackage{microtype}
\usepackage{amsmath}
\usepackage{booktabs,tabularx,array}
\usepackage{enumitem}
\usepackage{csquotes}
\usepackage{xcolor}
\usepackage[hidelinks]{hyperref}
\usepackage{titlesec}
\usepackage{fancyhdr}

\definecolor{accent}{HTML}{1C5CAB}
\definecolor{muted}{HTML}{5E636D}
\definecolor{warn}{HTML}{B4561E}
\titleformat{\section}{\Large\bfseries\color{accent}}{\thesection}{0.8em}{}
\titlespacing*{\section}{0pt}{18pt}{8pt}
\setlength{\parindent}{0pt}
\setlength{\parskip}{4pt}
\pagestyle{fancy}\fancyhf{}\fancyfoot[C]{\small\color{muted}\thepage}\fancyhead[R]{\small\color{muted}Viva preparation}
\renewcommand{\headrulewidth}{0pt}

% the report's results and names
\newcommand{\setR}[2]{\expandafter\def\csname R@#1\endcsname{#2}}
\newcommand{\R}[1]{\ifcsname R@#1\endcsname\csname R@#1\endcsname\else\pending\fi}
\newcommand{\pending}{\textit{\color{warn}(pending)}}
\InputIfFileExists{../report/generated/numbers.tex}{}{}
\newcommand{\drtc}{DRTC}
\newcommand{\smolvla}{SmolVLA}
\newcommand{\groot}{GR00T~N1.5}
\newcommand{\pizerofive}{$\pi_{0.5}$}
\newcommand{\degree}{$^\circ$}

% one question: number, question, answer; \push = what to add if the examiner pushes; \trap = the mistake to avoid
\newcounter{qn}
\newcommand{\q}[2]{\stepcounter{qn}\par\medskip
  {\bfseries\color{accent}Q\theqn.\ #1}\par\nopagebreak[4]#2\par}
\newcommand{\push}[1]{\par{\small\color{muted}\textbf{If pushed:} #1}}
\newcommand{\trap}[1]{\par{\small\color{warn}\textbf{Avoid:} #1}}
\newcommand{\todo}[1]{{\color{warn}[#1]}}

\begin{document}

{\LARGE\bfseries Viva preparation}\par\smallskip
{\large Reliable remote inference for low-cost robot manipulation under network faults}\par\smallskip
{\color{muted}Likely questions with short answers from our own results. Every number is taken from the report's
generated numbers, so it matches the report. Orange \pending{} marks results still running (setup v2 E1 + E2, E4).}

\tableofcontents

% ------------------------------------------------------------------------------------------------------------------
\section{The two-minute answer}

\textbf{What did you do?} Large robot policies (vision-language-action models of 0.45--3.4 billion parameters) are
too big for the computer on a low-cost arm, so they run on a GPU server and the robot sends its camera images over
the network for every chunk of actions. We built the whole system on LeRobot~0.7 for an SO-101 arm: datasets, seven
trained policies, remote inference, network fault injection and an evaluation framework. We then asked what network
faults do to such a robot and how to make it robust.

\textbf{What did you find?} On the robot, a mix of a 1~s latency spike, 1~s of lost replies and a 1.5~s outage took
the task success of LeRobot's shipped client from \R{main/smolvla/baseline/M1} to \R{main/smolvla/faults/M1}
(\smolvla) and from \R{main/groot/baseline/M1} to \R{main/groot/faults/M1} (\groot): the client ended every faulted
trial itself, one second after its action queue ran dry. With our implementation of Distributed Real-Time Chunking
(\drtc), with four fixes we found on the real setup, success under the same faults was \R{main/smolvla/drtc/M1} and
\R{main/groot/drtc/M1}, and the arm moved again within about 1.5~s of every fault. Smoothing the commands reduced the
jerk \R{main/smolvla/jerkratio}$\times$ and \R{main/groot/jerkratio}$\times$ without lowering success.

\textbf{Why should I believe it?} Because the first run's baseline gave up by design, we redesigned the experiments
(setup~v2): \drtc{} against a \emph{patient} client with the same limits, ablations that take \drtc{} apart, each
smoothing method alone before combining them, two tasks and seven policies on a replay bench, every plan frozen and
every comparison named before the data.

\subsection*{Numbers to have ready}
\begin{tabularx}{\textwidth}{X l}
\toprule
Round trip on the laptop: \smolvla{} / \groot{} & \R{main/smolvla/baseline/M4med} / \R{main/groot/baseline/M4med} \\
\quad of which inference on the A100 & \R{main/smolvla/baseline/M6} / \R{main/groot/baseline/M6} \\
\quad network share (two uncompressed 800$\times$600 images, 2.9~MB) & about 260~ms \\
First run: trials & \R{main/n} (2 policies $\times$ 4 conditions $\times$ 10) \\
Success \smolvla: baseline, faults, \drtc, \drtc{} + smoothing & \R{main/smolvla/baseline/M1}, \R{main/smolvla/faults/M1},
  \R{main/smolvla/drtc/M1}, \R{main/smolvla/smooth/M1} \\
Success \groot: same order & \R{main/groot/baseline/M1}, \R{main/groot/faults/M1}, \R{main/groot/drtc/M1},
  \R{main/groot/smooth/M1} \\
Faults vs \drtc{} (Fisher): \smolvla{} / \groot & \R{main/smolvla/p/faults-drtc/M1} / \R{main/groot/p/faults-drtc/M1} \\
Jerk reduction by smoothing & \R{main/smolvla/jerkratio}$\times$ / \R{main/groot/jerkratio}$\times$ \\
\groot{} stall fraction: \drtc{} alone / with smoothing & \R{main/groot/drtc/M2} / \R{main/groot/smooth/M2} \\
Raspberry Pi~5 as robot computer: loop rate; round trips & \R{edge/hz}; \R{edge/smolvla/rtt} / \R{edge/groot/rtt} \\
Bench network & 250 $\pm$ 20~ms, 8 runs per cell, 14~s per run \\
\bottomrule
\end{tabularx}

% ------------------------------------------------------------------------------------------------------------------
\section{Motivation and the problem}

\q{Why run the policy remotely at all?}{The policies worth using are vision-language-action models with 0.45 to 3.4
billion parameters. A laptop GPU can serve the smallest (\smolvla), but a Raspberry Pi or a low-cost robot computer
cannot run any of them at 30 actions per second. The usual answer, also LeRobot's, is a policy server on a GPU
machine; the price is that every action now depends on the network.}

\q{Why is the network a problem if the round trip is only 0.3--0.45~s?}{Because the robot plays the chunk while it
waits for the next one. If the next chunk is late, the queue runs dry and the arm stops (a stall). For \groot{} the
chunk is 16 actions, 0.53~s, barely longer than its 0.45~s round trip, so even a normal network leaves little margin.
And real networks are not steady: WiFi has spikes, lost packets and short outages.}

\q{What exactly goes wrong with LeRobot's shipped client?}{Two give-up limits. After 1~s with an empty queue it ends
the run with a fault, and when the queue runs dry it forgets the request in flight, so the late reply that would have
saved the run is thrown away. Also a request unanswered for 5~s ends the run. It keeps one request in flight and never
re-sends a lost one. In the first run it ended all 20 faulted trials at the first fault; on the bench each fault type
alone was enough (\R{bench/spike/smolvla/faults/M1f} \smolvla{} runs ended by the spike alone).}
\push{The policies were fine: they succeeded in most trials on the normal network. The failures came from the
client's rules.}

\q{Isn't the fix just to raise the timeout?}{That removes the give-up but not the weakness. A patient client still
keeps only half a second of actions in reserve and never replaces a lost reply: after the dropped replies it waits,
holding the arm, until the end of the run. In our replay checks a patient client stalled for over half of a faulted
run, with one stall of 7.5~s. That is exactly why setup~v2 compares \drtc{} with the patient client, not the shipped one.}
\trap{Do not claim the first run proves \drtc's scheduling is better than waiting; it proves the default is
fragile. Setup~v2 E1 answers the scheduling question: \R{v2/sort_one/e1/n} bench runs on sorting.}

% ------------------------------------------------------------------------------------------------------------------
\section{\drtc{} and our contribution}

\q{What is \drtc?}{Distributed Real-Time Chunking (Jack Vial) changes \emph{when} the robot asks for actions and
\emph{what it does} when a reply is late. Three parts: (1) a trigger that asks early, when the queue falls to
$H-s_{\min}$ actions; (2) a latency estimate in the style of TCP (Jacobson--Karels: smoothed mean plus $1.5\times$ the
smoothed deviation); (3) a cooldown re-request: a request unanswered within its cooldown is given up and a fresh
observation is sent; the freshest reply always wins.}

\q{What is yours and what is Vial's?}{The idea and the three mechanisms are Vial's. Ours: (a) the implementation
inside LeRobot~0.7's own runtime, replacing only the decision of when to request and when to give up, so all of
LeRobot's safety checks stay; (b) four fixes found on the real setup (next question); (c) six smoothing methods and a
rule to combine them; (d) the fault injector, the evaluation framework and all the experiments on seven policies.}

\q{What were the four fixes?}{All appeared only with a slow link and short chunks.
\begin{enumerate}[nosep]
\item \groot{} never received a chunk: the cooldown was capped at half a chunk, 8 steps = 0.27~s, while replies took
  0.4--0.7~s, so every request was given up before its answer. Fix: the cooldown is never shorter than the measured
  latency (the old cap is ablation A4).
\item A storm of re-requests: replies near the estimate were given up, the server, still busy, slowed further. Fix:
  the cooldown also covers the slowest of the last 10 replies.
\item The robot never started on a slow link: before the first reply the cooldown was only $s_{\min}+\varepsilon$
  steps. Fix: wait up to the 10~s start-up timeout for the first reply.
\item \enquote{Invalid RTC delay} errors: the requested inference delay could exceed the server's chunk. Fix: cap it.
\end{enumerate}}

\q{Why Jacobson--Karels and not just the maximum of recent replies?}{Upstream uses the maximum of the last 100
replies; one slow reply then inflates the estimate for 100 requests. Jacobson--Karels tracks the mean and the spread
and recovers in a few replies, the same reason TCP uses it. We kept a guard on the slowest of the last 10 for the
cooldown (fix 2), so we get a responsive estimate without giving up on replies that are merely slow. Ablation A2
replaces our estimator with upstream's, and on sorting it changed the stall by a median of \R{v2/sort_one/e2/a2_upstream_estimator/mix/delta} points under the fault mix: on the bench's steady link the two estimates are nearly the same. Be honest that the estimator is the part that did not matter here; it should matter on a link whose latency drifts.}

\q{How did you choose $s_{\min}$, $k$, the margins?}{$s_{\min}=20$ steps for 50-action policies and 8 for 16- and
24-action ones, so the request leaves about one round trip of actions; $k=1.5$ and $\alpha=1/8$, $\beta=1/4$ as in
TCP (RFC~6298). They were not tuned on the test data. Ablation A5 runs $s_{\min}$ at half and one and a half times
its value to show the sensitivity.}

\q{Does \drtc{} make the robot unsafe, e.g.\ acting on old observations?}{No more than upstream: LeRobot's checks stay
in place (observation freshness, execution generations, deadlines, latched faults), and the freshest observation
always wins, so after a fault the first chunk is computed from the current scene. The policies are closed-loop and
continue the task from wherever the arm stopped.}

% ------------------------------------------------------------------------------------------------------------------
\section{Smoothing}

\q{Where does jerky motion come from?}{Four sources: \emph{seams} (a new chunk from a newer observation does not
continue where the old one was heading), \emph{stalls} (the arm stops and restarts), \emph{jitter} (the policy's own
noise inside a chunk) and \emph{fault restarts} (after an outage the first new action can be far from the held pose).
Each method targets one source.}

\q{Explain each method in one line.}{\begin{itemize}[nosep]
\item Butterworth filter: causal 2nd-order low-pass at 3~Hz on every command (jitter, seams); costs lag.
\item Zero-lag smoothing: when a chunk arrives, filter the recent past and the queued future forward and backward
  (zero phase), so no lag (jitter).
\item RTC: the server is told how many actions will run during inference and guides the flow-matching denoising to
  continue them (seams at the source); only flow-matching policies with absolute actions (\smolvla, \groot).
\item Temporal ensembling: each step is a weighted average of every chunk that predicted it, as in ACT (seams).
\item Adaptive playback: when the queue would run dry before the next chunk, play slower (down to half speed) instead
  of stopping (stalls).
\item Resume ramp: after a stall, blend from the held pose into the stream with a minimum-jerk profile (restarts).
\end{itemize}}

\q{How did you measure smoothness?}{Jerk (M7): the RMS of the third difference of the commanded arm joints at 30~Hz,
in $10^3\,$\degree/s$^3$, \emph{only while the arm moves}. Counting idle time would make a stalling system look smooth,
which the first version of our metric did. Also the largest single step (M8b) and the disagreement at a chunk switch
(M8a).}

\q{Couldn't smoothing just lag behind the policy and look smooth?}{Yes, which is why we added M12, the deviation from
the plan: on every tick, the distance between what the arm was sent and what the policy asked for. The stack rule
refuses a method that adds more than 1\degree{} of it, because the bench has no task success to show the cost of lag.}

\q{Which smoothing worked? (setup v2, E3a)}{Each method did one job, the same on both tasks. Butterworth cut the jerk
to 0.12--0.28 of \drtc{} alone for every policy, at 0.9--1.6\degree{} deviation, without changing stalls. Adaptive
playback cut the stalls of the slow policies (\groot{} 64\,$\to$\,40\%, MolmoAct~2 55--57\,$\to$\,30--35\%, Diffusion
48--53\,$\to$\,29--30\%) with zero deviation. Combined they gave both. The ramp helped where the arm restarts often
(\groot's jerk to 0.18). Zero-lag helped the fast policies only; ensembling changed little.}

\q{Why didn't RTC help?}{It made each request slower: \smolvla's stalls more than doubled (about 14\,$\to$\,36--38\%)
and the smoother seams did not outweigh it. For \groot{} the jerk change was within the run-to-run spread. In the
stack test (E3b) the stack \emph{without} RTC was better than with it on all four pairs: \smolvla's stalls stayed at
\drtc's level (13--15\% instead of 31--33\%), and \groot's recovery from the latency spike on cube in box was 0.7~s
instead of 2.3~s.}
\push{RTC is designed for the case where inference is long relative to the chunk; here the network dominates the
round trip, so delaying inference further costs more than the seam saves. The first RTC request also compiles the
guided path and needs a warm-up.}

\q{How was the final stack chosen?}{By a rule written before the bench ran, applied by a script: the stack is adaptive
playback + resume ramp + one seam method (RTC or ensembling) + one jitter method (zero-lag or Butterworth). Per policy
and task, a candidate is ruled out if it adds over 1\degree{} of deviation, or raises stalls or any recovery time by
over 10\%, or recovers in fewer runs; of the rest, the lower relative jerk wins the pair; the method winning more pairs
is chosen. Result: \R{v2/stack}.}
\push{Both wins were thin: RTC won one pair and ensembling none; Butterworth won two of fourteen, zero-lag one; most
pairs ruled out both. We report that honestly, and E3b shows RTC is the weak part. \todo{state here what was decided
for E4: stack as frozen, or without RTC}}
\trap{Do not hide that the rule chose RTC and E3b then showed it hurts. Pre-registration means the rule stands; the
ablation is the evidence that corrects it.}

% ------------------------------------------------------------------------------------------------------------------
\section{Evaluation design and statistics}

\q{Why both a robot and a bench?}{The robot gives task success but costs operator time, so only a few conditions and
ten trials each. The bench (a replay robot that plays held-out recorded episodes as camera and joints, the same
servers, the latency measured on the robot simulated on every reply) runs thousands of trials under every fault type,
but nothing moves, so it has no task success. The robot answers \enquote{does the task succeed}, the bench
\enquote{what does each fault do to timing and motion}.}

\q{How is the bench realistic?}{Same code path: the same robot client, scheduler, smoothing and LeRobot policy server;
only the camera and arm are replaced by recorded held-out episodes, and the network by a simulated 250~$\pm$~20~ms
link, the network share measured on the robot. It reproduced the robot's failures (the shipped client ended 7--8 of 8
runs under each single fault) and found three of our four fixes in minutes. A run whose loop drops below 27~Hz is not
counted.}

\q{What is a frozen plan, and why?}{Each study is a plan file (models, conditions, trials, success rule, time limit,
comparisons, which are primary). Freezing pins every model to its exact Hub commit; servers only serve that commit.
So a result can only be made one way and can always be traced to its raw trials, and no model can change mid-study.
Plans may only grow (more trials, more models); anything else is a new plan.}

\q{How do you avoid order effects on the robot?}{Rounds: each round runs every model under every condition once, in
an order seeded by the plan's seed and the round number; layouts rotate across rounds. Drifts (light, motor
temperature, operator fatigue) hit all conditions alike.}

\q{Who judged success, and isn't that biased?}{The operator, by a success rule written in the plan before the first
trial; the operator knows the condition. A trial ended by the system is labelled \texttt{fault} from the run's own
record, whatever key was pressed. A setup mistake can void a trial with a reason, but voids are kept and listed, and
never used for a policy failure. Remaining bias is a stated limitation.}

\q{Why the fault shift?}{With fixed times every fault would hit the reach. Each trial's faults move by 0--4~s, seeded
by the trial number only, so trial $k$ meets its faults at the same moments in every condition and policy (comparisons
stay paired by trial) while faults also hit the grasp and the place.}

\q{Which statistics, and why?}{Success: Wilson 95\% intervals (good at small $n$ and near 0 or 100\%); differences in
success: Fisher's exact test. Other metrics: one value per trial, median and IQR, two-sided Mann--Whitney $U$ (no
normality assumption; heavy tails). Comparisons only within one plan and the same policy.}

\q{With so many conditions, aren't some results significant by chance?}{Yes, so each plan names its comparisons in
advance and a few as primary; only those are tested, with the Holm correction over the plan's primary tests. Everything
else is reported as exploratory with effect sizes and uncorrected $p$-values.}

\q{Is ten (or eight) trials enough?}{For large effects only. Ten robot trials give wide intervals (\smolvla{} with
\drtc: \R{main/smolvla/drtc/M1} with interval \R{main/smolvla/drtc/M1ci}). The shipped-client collapse (0 of 20) is
far outside that, but smoothing's success gain is not significant. On the bench, eight runs per cell: the same cell of
\groot{} gave jerk 178 in one plan and 219 in another, so we treat differences below about 20\% as noise. Robot plans
can grow to twenty trials.}

\q{Why is stall fraction a fair primary metric for E1 instead of success?}{The bench has no success, and against a
patient client the question is no longer whether the run is ended but how long the arm stands still and how quickly
it moves again after each fault (M11 recovery). Those are what \drtc's mechanisms change.}

% ------------------------------------------------------------------------------------------------------------------
\section{Results}

\q{Summarise the first run.}{Sorting task, laptop + A100, 80 trials. Faults took the shipped client to
\R{main/smolvla/faults/M1} and \R{main/groot/faults/M1} (it ended all 20 faulted trials). \drtc{} under the same faults:
\R{main/smolvla/drtc/M1} and \R{main/groot/drtc/M1} (\R{main/smolvla/p/faults-drtc/M1} and
\R{main/groot/p/faults-drtc/M1} against the faulted baseline), not significantly different from the fault-free
baseline, every fault recovered. With smoothing \R{main/smolvla/smooth/M1} and \R{main/groot/smooth/M1}, jerk
\R{main/smolvla/jerkratio}$\times$ and \R{main/groot/jerkratio}$\times$ lower.}

\q{Why is \smolvla{} with \drtc{} (\R{main/smolvla/drtc/M1}) below its baseline (\R{main/smolvla/baseline/M1})?}{The
interval is wide (\R{main/smolvla/drtc/M1ci}); the difference from the baseline is not significant. Faults still
cost something: a stall in the middle of a grasp can move the cube. Smoothing brought it to
\R{main/smolvla/smooth/M1}.}

\q{Why does \groot{} stall so much even without faults?}{Its 16-action chunk (0.53~s) is barely longer than its
0.45~s round trip, so the queue empties before the next chunk is ready: stall fraction \R{main/groot/baseline/M2} on
the normal network. This is chunk length versus latency, not the faults. Adaptive playback cut it to
\R{main/groot/smooth/M2}, at the cost of a slightly slower trial (\R{main/groot/drtc/M1t} vs \R{main/groot/smooth/M1t}).}

\q{What did setup v2 show about \drtc{} against a patient client? (E1)}{On sorting, 2{,}912 bench runs: no
difference without faults; under the fault mix \drtc{} cut the stall from \R{v2/sort_one/e1/mix/act/patient} to \R{v2/sort_one/e1/mix/act/drtc}\,\% for ACT and \R{v2/sort_one/e1/mix/smolvla/patient} to \R{v2/sort_one/e1/mix/smolvla/drtc}\,\% for \smolvla{}, and recovered from every fault for six of seven policies, while the patient client recovered \R{v2/sort_one/e1/mix/smolvla/patientrec} times (\smolvla). After an outage the patient client mostly never recovered (\R{v2/sort_one/e1/outage/smolvla/patientrec}): it re-sends once when the queue runs dry, and if that request is lost too it waits for the 10\,s startup deadline. \push{The honest negative: under a rising latency \drtc{} stalled more for every policy (\smolvla{} \R{v2/sort_one/e1/rise/smolvla/patient} vs \R{v2/sort_one/e1/rise/smolvla/drtc}\,\%), because the cooldown gives up replies that are only slow.} Cube in box is \todo{fill when bench\_v2\_drtc\_cubebox\_v1 finishes}.}

\q{Which part of \drtc{} matters most? (E2)}{The cooldown, as expected before the data. Without it the stall rose
by a median of \R{v2/sort_one/e2/a1_no_cooldown/mix/delta} points under the fault mix (\R{v2/sort_one/e2/a1_no_cooldown/mix/worse} policies), \R{v2/sort_one/e2/a1_no_cooldown/outage/delta} under outages. The estimator made no difference (\R{v2/sort_one/e2/a2_upstream_estimator/mix/delta}), the early trigger a little (\R{v2/sort_one/e2/a3_upstream_trigger/mix/delta}). \drtc's original $H/2$ cap raised the stall by \R{v2/sort_one/e2/a4_half_chunk_cap/none/delta} points even without faults: our first fix holds for all seven policies. The same cooldown is what hurts under a rising latency (\R{v2/sort_one/e2/a1_no_cooldown/rise/delta} points without it): mention it yourself.}

\q{After a 1--1.5~s fault, the next chunk comes from an old observation. Is it not safer to pause than to execute
stale actions?}{The cooldown sends a \emph{fresh} observation, so the first chunk after a fault is planned from the
current scene; the aligned merge skips the steps the arm already executed while the request was in flight; and LeRobot
refuses any chunk whose observation is older than 5~s. What \drtc{} never does is execute an old chunk after a long
stall: it holds instead.}

\q{Your simulator and bench measure stalls. Does less stalling mean more task success?}{Not necessarily, which is why
the robot experiments measure success: in the first robot run \drtc{} brought success under faults back to the
fault-free level for both policies. The bench and the simulator explain the mechanism; success is judged on the arm
(E4), and possibly in simulation (LIBERO, proposed).}

\q{When the cooldown fires, what happens to the request still running on the GPU?}{It keeps running: a forward pass
cannot be cancelled. The replacement gets a BUSY reply, backs off one step and is sent again when the server is free,
so a cooldown that fires too early wastes server time. That is exactly why the cooldown is never shorter than the
measured latency and covers the slowest of the last ten replies (fixes 1--2), and why a rising latency still hurts.}

\q{What did the final stack do? (E3b)}{Against \drtc{} alone it lowered the jerk 3--20$\times$ and the largest step
2--6$\times$ for every policy, and cut the stalls of the slow policies (Diffusion 50\,$\to$\,27\%, MolmoAct~2
57\,$\to$\,32\%, GR00T~N1.7 68\,$\to$\,49\% on sorting). Costs: 0.7--2.3\degree{} deviation from the plan, more stalls for
\smolvla{} (because of RTC) and \pizerofive, and slower recovery for \pizerofive{} on sorting (1.0 to 2.2~s). Removing a part:
without adaptive playback the stalls come back; without Butterworth the jerk is about twice as high; the ramp matters little;
removing RTC lowered the stalls (\smolvla{} by 16--20 points, \groot{} by 3) at no cost in jerk.}

\q{What did the Raspberry Pi test show?}{The Pi~5 holds the 30~Hz loop (\R{edge/hz}, no dropped frames), so the
robot side fits on it. But the round trip roughly doubles (\R{edge/smolvla/rtt} and \R{edge/groot/rtt}) because each
request crosses the hotspot WiFi to the laptop first; with the shipped client \groot's queue ran dry even without
faults. The edge is exactly the regime \drtc{} and image compression are for; official Pi trials are future work.}

\q{Which policy is best?}{Depends on the criterion. Offline on sorting \smolvla{} had the lowest held-out arm error
(\R{offline/sortone/smolvla/M9arm}\degree), ACT was fastest (\R{offline/sortone/act/M10} per chunk) and smallest;
\groot{} was the largest (\R{offline/sortone/groot/M10gb}). For remote inference the chunk length relative to the
round trip decides how much it stalls: ACT and \smolvla{} barely stall, the 16-action GR00T models stall most.
\todo{add the robot baselines when run}}

% ------------------------------------------------------------------------------------------------------------------
\section{Limitations and threats to validity}

\q{What are the main limitations?}{\begin{itemize}[nosep]
\item Small robot samples (10 per cell): only large success differences are significant.
\item Robot \drtc{} experiments on one task and two policies; the bench covers two tasks and seven policies but has no
  task success.
\item The operator judges success and knows the condition (rule fixed in advance; fault labels from the record).
\item Faults have fixed shapes with a seeded shift, not a recorded real network; the bench simulates latency.
\item Shared GPU servers: inference times include other users' jobs.
\item Uncompressed images: 2.9~MB per request is most of the round trip.
\item Two trainings (\pizerofive, MolmoAct~2) stopped early for time.
\end{itemize}}

\q{Does the replay bench overstate or understate the problem?}{Its images do not react to the actions, so a policy
cannot correct itself on the bench; we never read task outcome from it, only timing and motion, which do not depend on
the scene reacting. The latency is the one measured on the robot. It reproduced the robot's failure pattern.}

\q{Two bench runs shared the GPU at one point. Does that bias results?}{The stack evals (V4) ran in parallel with the
\drtc{} evals on the same H200. Each eval runs one trial at a time and the GPU was idle between requests; the \drtc{}
eval's throughput did not change (3.35 runs per minute before and during). Contention affects all conditions of an
eval alike because trials are interleaved in rounds, and comparisons are only made within one eval.}

\q{Could the results be specific to your WiFi?}{The mechanisms are general (a queue drained by playback and refilled
over a link), and the bench sweeps the fault types rather than one network. The absolute numbers are specific; the
ranking of schedulers under each fault type is what we claim.}

\q{\pizerofive{} and MolmoAct~2 were not trained to the end. Is that a problem?}{For their task performance, possibly;
for this project's questions, little: the bench measures timing and motion, which depend on chunk length and inference
time, not on how well-trained the weights are. We kept the recipe's learning-rate schedule length so the stopped run is
identical to the full one up to that step, and say so in the report.}

% ------------------------------------------------------------------------------------------------------------------
\section{Implementation and engineering}

\q{What did you build besides \drtc?}{Recording and automatic trimming of datasets with a held-out split; one-command
training with each author's recipe, logged to Weights \& Biases and Trackio; a compatibility layer for older
checkpoints; the fault injector (drop, duplicate, reorder, spikes, latency windows, random loss, disconnects); the
evaluation framework (frozen plans, seeded rounds, raw records, a shared results store, reports rebuilt from raw files);
a Raspberry Pi setup; this report and the slides, built from the results. Over \R{software/tests} automated tests.}

\q{What was the hardest bug?}{\groot{} never receiving a chunk under \drtc: the cooldown cap at half a 16-action chunk
(0.27~s) was shorter than every round trip, so every request was given up and re-sent before its reply. It only shows
with a slow link and a short chunk, which is why the replay bench with the measured latency became our main
development tool.}

\q{A method that looked good but failed?}{The soft stop: slowing the last actions before the queue ran dry. Alone it
looked harmless, but by keeping the queue alive it let a reply from a 1.3~s old observation be accepted, and the next
fresh chunk jumped (largest step 38--42\degree{} against 18\degree). We dropped it, and from then on tested every
smoothing method alone before combining them.}

\q{How reproducible is this?}{Every plan is frozen to exact model commits; every trial keeps its raw records, which are
never edited; every table, figure and number in the report is rebuilt from them by one command; the training recipes are
the authors' own, fixed in the repository; the bench runs unattended with one script.}

% ------------------------------------------------------------------------------------------------------------------
\section{Bigger picture and future work}

\q{What would you do with three more months?}{Official edge trials on the Pi; compressing the images (most of the
round trip); RTC for relative-action policies such as GR00T~N1.7 (re-anchoring the leftover actions to the new state);
an on-robot fallback policy (ACT is small enough) that keeps the arm moving while the server is unreachable; robot
plans grown to twenty trials and to the cube-in-box task.}

\q{Would 5G or a wired link make this unnecessary?}{It would make faults rarer, not impossible, and a shorter round
trip still drains the queue during a fault. \drtc's mechanisms cost almost nothing on a good network (the no-fault
cells of E1 measure that cost), so they are cheap insurance.}

\q{Why not run a smaller policy on the robot and avoid the network?}{That trades capability for reliability: the
strongest policies on this arm in published SO-101 studies are the large VLAs (\pizerofive{} first). A small on-robot
fallback combined with remote inference is in our future work.}

\q{What is the one thing you want me to remember?}{Remote inference does not fail because the models are wrong but
because the client gives up or waits for a reply that never comes. Asking early, estimating the latency properly and
re-requesting with a fresh observation keeps the robot working through faults, and smoothing then makes the motion
smooth, as long as you measure what it costs in deviation from the plan.}

\end{document}
